% Resume -- Christopher Suarez
%
% Pushing a change to this file on main rebuilds assets/files/resume.pdf
% (see .github/workflows/build-resume.yml).
%
% Local preview (output lands next to this file and is git-ignored):
%   tectonic resume/resume.tex
% Needs XeLaTeX or LuaLaTeX (Tectonic is XeLaTeX); on Overleaf, set the
% compiler to XeLaTeX.

\documentclass[letterpaper,12pt]{article}

\usepackage{fontspec}
% Carlito is an open-source font metrically compatible with Calibri, the font
% the old Word version used, so text lines up the same way. Ligatures are off
% to match Word and so resume-parsing software reads every word cleanly.
\setmainfont{Carlito}[
  Extension      = .ttf,
  UprightFont    = *-Regular,
  BoldFont       = *-Bold,
  ItalicFont     = *-Italic,
  BoldItalicFont = *-BoldItalic,
  Ligatures      = {TeX, NoCommon},
  Numbers        = Lining,
]
% Body text sits 0.1in inside the 0.5in page margins; section headings and the
% header rule reach back out to the page margin (see \outdent).
\usepackage[top=0.75in,bottom=0.75in,left=0.6in,right=0.6in]{geometry}
\usepackage{xcolor}
\usepackage{enumitem}
\usepackage[hidelinks,pdfauthor={Christopher Suarez},pdftitle={Christopher Suarez Resume}]{hyperref}

\pagestyle{empty}
% Word's single spacing for Calibri is 1.22x the font size; LaTeX's is 1.2x.
\linespread{1.014}
\setlength{\parindent}{0pt}
\raggedright
\frenchspacing
\hyphenpenalty=10000

\newlength{\outdent}
\setlength{\outdent}{0.1in}

\definecolor{linkblue}{HTML}{0563C1}

% Underline at a fixed depth so letters with descenders don't push it lower.
\newcommand{\ul}[1]{\underline{\smash{#1}}}
% Hyperlink styled the way Word shows them: blue and underlined.
\newcommand{\link}[2]{\href{#1}{\textcolor{linkblue}{\ul{#2}}}}

% Bold, all-caps section title at the page margin.
\newcommand{\resumesection}[1]{%
  \par\vspace{4pt}%
  \hspace*{-\outdent}\textbf{#1}\par\nopagebreak}
% Organization or venue, with a right-aligned location or date.
\newcommand{\heading}[2]{%
  \par\vspace{2pt}%
  #1\hfill #2\par\nopagebreak}
% Italic job title, with right-aligned dates.
\newcommand{\subheading}[2]{\textit{#1}\hfill #2\par\nopagebreak}
% Italic degree, school, and right-aligned year.
\newcommand{\degree}[3]{\textit{#1}, #2\hfill #3\par}
% Bold project title within a job.
\newcommand{\project}[1]{\textbf{#1}\par\nopagebreak}
% Indented publication: linked label -- italic title.
\newcommand{\publication}[3]{\hspace*{0.2in}\link{#1}{#2} -- \textit{#3}\par}

\setlist[itemize]{label=o, nosep, leftmargin=0.4in, labelindent=0.2in,
  labelwidth=0.2in, labelsep=0pt, align=left, beginpenalty=10000}

\begin{document}

{\centering
  % The strut gives the name line Word's full ascent from the top margin.
  {\fontsize{20pt}{24pt}\selectfont \rule{0pt}{19pt}Christopher Suarez\par}
  \vspace{2pt}
  chriswsuarez@gmail.com \textbar{} 954-336-9953 \textbar{}
  \link{https://www.linkedin.com/in/chriswsuarez}{linkedin} \textbar{}
  \link{https://chriswsuarez.github.io/}{github}\par}
\vspace{-5pt}
\hspace*{-\outdent}\rule{\dimexpr\textwidth+2\outdent}{0.5pt}\hspace*{-\outdent}\par
\vspace{3pt}

\resumesection{SUMMARY}
Robotics engineer with 10 years experience focused on development \& deployment of autonomous systems

\resumesection{EDUCATION}
\degree{M.S. Mechanical Engineering}{University of Texas at Austin}{2019}
\degree{B.S. Mechanical Engineering}{University of Florida}{2014}

\resumesection{EMPLOYMENT}
\heading{\ul{Moderna}}{Boston, MA}
\subheading{Sr. Engineer II, Robotic Systems}{Nov. 2026 -- Present}

\heading{\ul{Apptronik}}{Austin, TX}
\subheading{Sr. Systems Integration Engineer}{May. 2026 -- Oct. 2026}
\begin{itemize}
  \item Redesigned, implemented, and deployed the reprovisioning process for Windows teleop machines
  \item Automated critical fleet-wide deployment and diagnostic processes via Ansible \& crontab on server
  \item Supported and improved internal fleet monitoring through Zabbix
  \item Implemented harbor docker registry mirror to mitigate local network bandwidth issues across fleet
\end{itemize}
\subheading{Sr. Software Support Engineer}{July 2025 -- May. 2026}
\begin{itemize}
  \item Primary software engineer responding to software / firmware issues with deployed robots
  \item Performed all OTA software upgrades to the fleet and PR/MR'd various OTA software improvements
  \item Supported critical robot deployment and \link{https://youtu.be/UObzWjPb6XM?si=ZYu0lCAWmv4WC2iq&t=219}{demo} events domestically and internationally
  \item Created diagnostic tools for Field Service and Operations teams for ECAT and firmware diagnosis
\end{itemize}

\heading{\ul{TORC Robotics (contractor through Actalent)}}{Austin, TX}
\subheading{Software Integration Engineer}{Nov. 2024 -- May. 2025}
\begin{itemize}
  \item Streamlined Jenkins CI/CD pipeline using Bazel, reducing CI build times by 33 -- 50\%
  \item Created CPU \& Memory resource usage tool to assess sensor driver performance at runtime and in CI
  \item Updated IMU and Lidar sensor drivers to modern internal fault framework
  \item Debugged hardware and sensor firmware issues for Hardware-in-the-loop (HIL) Lab
  \item Expanded sensor validators to new sensor drivers ensuring performance to meet requirements
\end{itemize}

\heading{\ul{UT Austin Nuclear \& Applied Robotics Group}}{Austin, TX}
\subheading{Research Engineering Scientist Associate IV}{Sept. 2019 -- Nov. 2024}
\project{Summary}
\begin{itemize}
  \item Oversaw and performed over 100 technical demos directly leading to \textgreater\$10M in project funding to UT
  \item Expertise in setting up, debugging, and operating robotics systems with ROS
  \item Managed development teams for multiple research integration projects
  \item Mentored new engineer and student hires across varying research efforts
\end{itemize}
\project{Army Futures Command UTDD Project SPEED}
\begin{itemize}
  \item Lead demo engineer overseeing a large-scale on-site technical demonstration with a team of 12
  \item Stood up a new applied research lab focusing on large vehicle outdoor autonomy
  \item Integrated Robot Technology Kernel (RTK) 22 onto a previously incompatible Hunter WOLF platform
  \item Created a simulation test environment in Unreal Engine
\end{itemize}
\project{Augmented Robot Environment (\link{https://utnuclearroboticspublic.github.io/Augmented-Robot-Environment/}{AugRE}) Project}
\begin{itemize}
  \item Developed and dockerized software stack to co-localize robots with Hololens2 AR Headsets
  \item Managed dev team creating software for AR in Unreal Engine and robotics in ROS/C++
  \item Setup and managed LAN \& WAN network for multiple autonomous and human agent comms
  \item Performed dozens of live hands-on demonstrations to VIPs, Gov Officials, Army Generals, CEOs
\end{itemize}
\project{Army Futures Command UTDD Integrated Prototype Project}
\begin{itemize}
  \item Contributed features, sensor integrations, and critical fixes to Army Research Labs' autonomy stack
  \item Integrated ARL Phoenix navigation stack with custom research platform hardware
  \item Conducted multi-week study to quantify long-term \link{https://www.youtube.com/watch?v=QU-jWwfyaAE}{autonomy} viability of various UT research efforts
  \item Created custom robot and environment assets for Unity environment for simulation development
  \item Developed object persistent localization \link{https://github.com/UTNuclearRobotics/object_detection/tree/main}{software} fusing camera, LIDAR, and mobile robot data
\end{itemize}
\project{AI2C ATR-MCAS Project}
\begin{itemize}
  \item Stood up an Army research platform for autonomous navigation using GPS, IMU, LIDAR, Segway base
  \item Incorporated capability to visualize our autonomous software stack with ATAK (Android Tactical Awareness Kit) via Cursor-On-Target (CoT) XML messaging
  \item Extended ATAK visualization capability to our AR AugRE system allowing human agents to appear on ATAK and interface with autonomous agents without co-localizing directly with AR anchors
  \item Performed yearly deliverable hands-on live demonstrations on-site for Army VIPs
\end{itemize}
\project{SRNL H-Canyon Tunnel Inspection Project}
\begin{itemize}
  \item Developed 3D obstacle detection software for an off-road dual arm mobile manipulator via ROS / C++
  \item Designed and constructed a 9'x12'x36' above ground tunnel as a mockup harsh test environment
  \item Performed final semi-autonomous \link{https://www.youtube.com/watch?v=X-VfCXZ3_cE&t=3s}{demonstration} of tunnel inspection inside a mock environment
\end{itemize}

\heading{\ul{Lockheed Martin Missiles \& Fire Control}}{Orlando, FL}
\subheading{Systems Engineer}{June 2014 -- July 2016}
\begin{itemize}
  \item Created system level engineering requirements and test procedures
  \item Conducted program critical tests in-plant and on-site - domestically and internationally
  \item Designed GUIs to operate and test complex machinery and contributed to the end product UI design
\end{itemize}

\resumesection{PERSONAL PROJECT}
\heading{\link{https://baseballsimpleton.com/}{Baseball Simpleton}}{Apr. 2026}
\begin{itemize}
  \item Analyzed MLB data from statcast era into accessible UI (Typescript / React) for player analysis
  \item Split data into year-by-year SQL databases and one meta SQL database to allow for fast queries
  \item Implemented deployment through AWS EC2 system, DBs on S3 bucket, and crontab to update nightly
\end{itemize}

\resumesection{PUBLICATIONS \& PRESENTATIONS}
\heading{\ul{IROS 2023 XR-ROB Workshop}}{Oct. 2023}
\publication{https://arxiv.org/abs/2401.05665}{Abstract/Presentation}{AR UI for Command, Control, and Supervision of Large Multi-Agent Teams}
\heading{\ul{IEEE International Conference on Robot \& Human Interactive Communication (RO-MAN) 2022}}{Sept. 2022}
\publication{https://ieeexplore.ieee.org/document/9900721}{Paper}{AugRE: Augmented Robot Environment to Facilitate Human-Robot Teaming and Communication}
\heading{\ul{Waste Management Symposia 2019}}{Mar. 2019}
\publication{https://www.osti.gov/biblio/23005359}{Paper}{Savannah River Site H-Canyon Tunnel Inspection LiDAR Mapping Solution}
\publication{https://www.osti.gov/biblio/1501414}{Paper}{Savannah River Site H-Canyon Advanced Technology Demonstration (ATD) Program Review}

\resumesection{SKILLS}
\begin{itemize}
  \item Programming Languages: Python, C++, Typescript, Shell
  \item AI Tools: Claude Code, OpenAI Codex, Gemini
  \item Dev Tools: Git, Docker, Ansible, Zabbix, Bazel, Jenkins, AWS EC2 \& S3, Jira
  \item Simulation: Unreal Engine, Gazebo, Unity
  \item Software: ROS, Solidworks, Visual Studio, MatLAB, LabVIEW, QT, Linux, MacOS
  \item IT: LAN / WAN networking, embedded machines, edge compute
  \item Hardware: Additive manufacturing, lathes, mills, CNC\ldots
  \item Systems: Requirements, validation \& testing, technical reporting, multi-disciplinary collaboration
\end{itemize}

\end{document}
